\documentclass[12pt]{article}
\usepackage{fontspec}
\usepackage{polyglossia}
\setmainlanguage{french}
\usepackage{csquotes}

\usepackage{tikz}% faire des schémas
\usepackage{pgfplots}% graphiques à partir de données (courbes, histogrammes)
\usepackage{pgfplotstable}% importer les données depuis un fichier csv
\pgfplotsset{compat=1.16}
\pgfplotsset{/pgf/number format/1000 sep=}% pas de séparateur pour les nombres au-dessus de mille
% 4. Le fichier de l'Insee utilise la virgule décimale (0,9) : elle doit être lue comme un point.
\pgfplotsset{/pgf/number format/read comma as period}
% 5. Réglage de tous les tableaux importés : titres de colonnes en gras, colonnes alignées à gauche.
\pgfplotstableset{%
	assign column name/.style={/pgfplots/table/column name={\textbf{#1}}},
	column type=l}

% Les fichiers de données sont dans le dossier exercices : d'où le chemin ../exercices/
% Le fichier nettoyé (question 4) est dans le dossier corrections.

\begin{document}

\section{Du fichier CSV au tableau}

% 1. \pgfplotstableread lit le fichier et stocke les données dans la variable \travail.
%    col sep=comma : les colonnes sont séparées par des virgules.
%    Attention : ne pas appeler la variable \table, nom déjà utilisé par l'environnement table.
%    \pgfplotstabletypeset affiche le contenu de la variable sous forme de tableau.

\pgfplotstableread[col sep=comma]{../exercices/travail_domestique_insee.csv}{\travail}

\pgfplotstabletypeset{\travail}


\section{Du fichier CSV au graphique}

% 2. Les mêmes données, sans rien recopier : \addplot table lit la variable.
%    x= et y= : noms des colonnes, tels qu'ils figurent dans la première ligne du fichier.

\begin{tikzpicture}
	\begin{axis}[xtick=data]
		\addplot table[x=Date, y=Hommes]{\travail};
		\addplot table[x=Date, y=Femmes]{\travail};
		\legend{Hommes,Femmes}
	\end{axis}
\end{tikzpicture}


\section{Des données textuelles}

% 3. Les prénoms ne sont pas des nombres : ils ne peuvent pas servir d'abscisse.
%    - x expr=\coordindex : l'abscisse est le numéro de la ligne (0, 1, 2...) ;
%    - xticklabels from table={\prenoms}{Prénom} : les étiquettes de l'axe sont lues dans la colonne Prénom ;
%    - x tick label style={rotate=90} : les étiquettes sont écrites verticalement.

\pgfplotstableread[col sep=comma]{../exercices/classement_prenoms_insee.csv}{\prenoms}

\begin{tikzpicture}
	\begin{axis}[ybar, xtick=data, xticklabels from table={\prenoms}{Prénom}, x tick label style={rotate=90}]
		\addplot table[x expr=\coordindex, y=Nombre]{\prenoms};
	\end{axis}
\end{tikzpicture}


\section{Nettoyer ses données}

% 4. Nettoyage dans le tableur (fichier insee_categories_socio_professionnelles.csv, dossier corrections) :
%    - suppression du titre, de l'unité (« en % »), des notes, de la source et des lignes vides ;
%    - suppression de l'appel de note dans « Employés1 » ;
%    - suppression des sous-catégories (« Employés qualifiés »...), déjà comptées dans « Employés » et « Ouvriers » ;
%    - enregistrement en CSV avec la tabulation comme séparateur,
%      car « Artisans, commerçants, chefs d'entreprise » contient des virgules.
%    La première colonne contient du texte : on le signale avec string type.

\pgfplotstableread[col sep=tab]{insee_categories_socio_professionnelles.csv}{\csp}

\pgfplotstabletypeset[columns/Catégorie socioprofessionnelle/.style={string type}]{\csp}

\end{document}
